\documentclass{article}

\usepackage{amsmath}
\usepackage{amssymb}

\newcommand{\grad}{\nabla}
\newcommand{\divergence}{\nabla\!\cdot}

\begin{document}

\section*{Immiscible two-phase flow without gravity}

Let $\Omega\subset\mathbb{R}^{d}$ be the porous domain and
$t\in(0,T]$.  The phases are water ($w$) and carbon dioxide ($c$).
The primary unknowns are their pressures
\[
    p_w,p_c:\Omega\times(0,T]\longrightarrow\mathbb{R}.
\]
There is no dissolved carbon dioxide, no mass transfer between phases,
and gravity is neglected.  The capillary pressure (suction) is
\[
    s = p_c-p_w,\qquad s\geq 0.
\]
The water and carbon-dioxide saturations satisfy
\[
    S_w=S_w(s),\qquad S_c=1-S_w,\qquad
    S_w,S_c\in[0,1].
\]

\subsection*{Mass conservation}

For each phase $\alpha\in\{w,c\}$, define its Darcy velocity by
\[
    \boldsymbol{u}_{\alpha}
    =
    -\frac{k_{r\alpha}}{\mu_\alpha}\,
     \boldsymbol{K}_{\mathrm{int}}\grad p_\alpha ,
\]
where $n$ is the porosity, $\boldsymbol{K}_{\mathrm{int}}$ is the
intrinsic-permeability tensor, $\mu_\alpha$ is the dynamic viscosity,
and $k_{r\alpha}$ is the relative permeability.  With no source terms,
the phase mass balances are
\[
    \frac{\partial}{\partial t}\left(nS_\alpha\rho_\alpha\right)
    +\divergence\left(\rho_\alpha\boldsymbol{u}_\alpha\right)=0.
\]
Thus, the coupled pressure system is
\begin{align}
    \frac{\partial}{\partial t}
    \left[
        nS_w(p_c-p_w)\rho_w(p_w)
    \right]
    -
    \divergence\left[
        \rho_w(p_w)
        \frac{k_{rw}(S_w(p_c-p_w))}{\mu_w}
        \boldsymbol{K}_{\mathrm{int}}\grad p_w
    \right]
    &=0,
    \label{eq:water_mass}
    \\
    \frac{\partial}{\partial t}
    \left[
        n\bigl(1-S_w(p_c-p_w)\bigr)\rho_c(p_c)
    \right]
    -
    \divergence\left[
        \rho_c(p_c)
        \frac{k_{rc}(S_w(p_c-p_w))}{\mu_c}
        \boldsymbol{K}_{\mathrm{int}}\grad p_c
    \right]
    &=0.
    \label{eq:co2_mass}
\end{align}

\subsection*{Constitutive laws}

The water-retention law is
\[
    S_w(s)
    =
    S_{\mathrm{res}}
    +(1-S_{\mathrm{res}})
    \left[
        1+\left(\frac{s}{P_r}\right)^m
    \right]^{\frac{1}{m}-1},
\]
where $S_{\mathrm{res}}\in[0,1)$ is the residual water saturation,
$P_r>0$ is a reference capillary pressure, and $m>0$ is a model
parameter.

The compressible phase densities are
\[
    \rho_w(p_w)
    =
    \rho_{w0}
    \exp\left[\kappa_T(p_w-p_{wr})\right],
    \qquad
    \rho_c(p_c)
    =
    \frac{M_c}{ZRT}\,p_c.
\]
Here $\rho_{w0}$ is a reference water density, $\kappa_T$ is the water
compressibility, $p_{wr}$ is the reference water pressure, $M_c$ is
the molar mass of carbon dioxide, $Z$ is its compressibility factor,
$R$ is the universal gas constant, and $T$ is the absolute
temperature.

For a given $S_w$, define the effective water saturation
\[
    S_e=\frac{S_w-S_{\mathrm{res}}}{1-S_{\mathrm{res}}}.
\]
The relative permeabilities are
\begin{align*}
    k_{rw}(S_w)
    &=
    \sqrt{S_w}
    \left\{
        1-
        \left[
            1-S_w^{1/C_{kw,1}}
        \right]^{C_{kw,1}}
    \right\}^{2},
    \\
    k_{rc}(S_w)
    &=
    (1-S_e)^{C_{ka,1}}
    \left(1-S_e^{C_{ka,2}}\right).
\end{align*}

Equations~\eqref{eq:water_mass}--\eqref{eq:co2_mass}, together with
these constitutive laws and appropriate initial and boundary
conditions, determine $p_w$ and $p_c$.


\subsection*{Initial-boundary value problem}

Let the boundary be decomposed into the disjoint left, right, and
impermeable parts
\[
    \partial\Omega
    =
    \overline{\Gamma_L\cup\Gamma_R\cup\Gamma_N}.
\]
Initially, the pressure fields are uniform:
\[
    p_w(\boldsymbol{x},0)=15\,\mathrm{MPa},
    \qquad
    p_c(\boldsymbol{x},0)
    =
    p_w(\boldsymbol{x},0)+0.1\,\mathrm{MPa}
    =
    15.1\,\mathrm{MPa},
    \qquad \boldsymbol{x}\in\Omega.
\]
The right boundary is held at the initial reservoir pressures,
\[
    p_w(\boldsymbol{x},t)=15\,\mathrm{MPa},
    \qquad
    p_c(\boldsymbol{x},t)=15.1\,\mathrm{MPa},
    \qquad
    (\boldsymbol{x},t)\in\Gamma_R\times(0,T].
\]
Carbon dioxide is injected through the left boundary by prescribing a
pressure history $p_{\mathrm{inj}}(t)$:
\begin{align*}
    p_c(\boldsymbol{x},t)=p_{\mathrm{inj}}(t),
    \qquad
    (\boldsymbol{x},t)&\in\Gamma_L\times(0,T],\\
    p_{\mathrm{inj}}(0)&=15.1\,\mathrm{MPa},
    &
    p_{\mathrm{inj}}(t_1)&=20\,\mathrm{MPa}.
\end{align*}
For example, a linear pressure ramp followed by a hold is
\[
    p_{\mathrm{inj}}(t)
    =
    15.1\,\mathrm{MPa}
    +
    4.9\,\mathrm{MPa}\,
    \min\left(\frac{t}{t_1},1\right).
\]
No water crosses the injection boundary, and the remaining boundary is
impermeable to both phases:
\begin{align*}
    \rho_w\boldsymbol{u}_w\cdot\boldsymbol{n}&=0
    &&\text{on }\Gamma_L\times(0,T],\\
    \rho_\alpha\boldsymbol{u}_\alpha\cdot\boldsymbol{n}&=0,
    \quad \alpha\in\{w,c\},
    &&\text{on }\Gamma_N\times(0,T].
\end{align*}


\end{document}
